\chapter{Reading the Periodic Table}

\begin{kaichang}
This appendix is a tool rather than a story, like a multiplication chart beside your pencil sharpener. Treat it as the periodic table's instruction manual. Whenever the main text raises questions such as ``Which group is this element in?'', ``Why is fluorine fiercer than chlorine?'', or ``What does +2 valence mean?'', return here.

If you have read Chapters 10 and 11, you have seen most of this in story form. Here it is compressed into reference tables, coordinates, and rules. If you have not reached those chapters, no matter: each rule points you toward its fuller explanation.
\end{kaichang}

\section{First, Read One Square}

Each periodic-table square holds an element's identity card. Layout varies, but there are three core pieces of information:

\begin{itemize}[itemsep=2pt]
\item \textbf{Atomic number}: the integer in a corner, equal to nuclear proton count. This is the element's true number: hydrogen 1, carbon 6, oxygen 8, iron 26, gold 79. The whole table follows it in ascending order without gaps or duplicates. Chapter 7 explains why proton count defines identity.
\item \textbf{Element symbol}: one or two letters in the center, the first uppercase and the second lowercase. $\chem{Co}$ is cobalt; $\chem{CO}$ is carbon monoxide. A tiny capitalization difference means a completely different substance. Appendix A gives the notation rules.
\item \textbf{Relative atomic mass}: the longer number below, an average mass of the element's naturally mixed isotopes, conventionally written without units. Chlorine's 35.45 is not the mass of a particular atom but a weighted average of two isotopes.
\end{itemize}

Many tables color the squares by metals, nonmetals, and metalloids. These colors are human labels, not actual appearances. Gold is yellow, but most metals are silvery gray.

Practice with three squares. Find number 11: $\chem{Na}$, sodium, whose positive ion contributes table salt's salty taste. Number 17 is $\chem{Cl}$, chlorine, central to the disinfectant smell when pool water catches in your throat. Number 26 is $\chem{Fe}$, iron, involved both in red blood and reddish-brown rust. Three squares, three familiar scenes: the periodic table is never farther from daily life than your kitchen.

\section{Coordinates: Rows and Columns}

To use the table as a map, learn its coordinates.

\textbf{Horizontal rows are periods}, numbered 1 through 7 from top to bottom. Across a period, electrons are added one by one on the same ``floor,'' and behavior changes steadily from metals eager to lose electrons to nonmetals desperate to gain them. Period 1 has 2 squares; periods 2 and 3 have 8 each; periods 4 and 5 have 18 each; periods 6 and 7 have 32 each when the lanthanides and actinides are restored. These numbers directly reflect electron-floor capacities (Chapters 9 and 10).

\textbf{Vertical columns are groups}, eighteen altogether, numbered internationally from 1 to 18 left to right. You may also encounter older numbering: A for main groups, IA, IIA, through VIIIA, and B for the others. Both systems remain common; this book uses 1--18, so cross-reference older labels when needed.

Members of a group have the same outer-electron count and resemble a family. Group 1, excluding hydrogen, contains soft metals that explode on meeting water; group 17 excels at grabbing electrons; group 18 consists of noble gases that ignore everything. \textbf{To guess an unfamiliar element's behavior, inspect its upstairs and downstairs neighbors}. This is the table's most useful feature and how Mendeleev predicted unknown elements (Chapter 10).

A boundary slopes from boron, B, toward astatine, At. Metals lie below and left, nonmetals above and right, while metalloids such as boron, silicon, germanium, arsenic, antimony, and tellurium straddle it. Silicon's partly metallic and partly nonmetallic nature is exactly why it serves in your phone's chips.

\section{Four Blocks: Neighborhoods by Orbital Type}

A broader division than groups is \textbf{blocks}, classified by the type of room, an s, p, d, or f orbital, receiving the last electron:

\begin{table}[htbp]
\centering
\begin{tabular}{llll}
\toprule
Block & Position & Elements & Shared behavior \\
\midrule
s & \shortstack[l]{Leftmost 2 columns\\(groups 1, 2)} & \shortstack[l]{Hydrogen, alkali and\\alkaline-earth metals} & \shortstack[l]{1--2 outer electrons;\\easily lost; reactive} \\
p & \shortstack[l]{Rightmost 6 columns\\(groups 13--18)} & \shortstack[l]{Carbon, nitrogen, oxygen,\\halogens, noble gases, etc.} & \shortstack[l]{Metals and nonmetals;\\many life elements} \\
d & \shortstack[l]{Middle 10 columns\\(groups 3--12)} & \shortstack[l]{Transition metals: iron,\\copper, zinc, titanium, etc.} & \shortstack[l]{Multiple valences;\\colored ions; catalysts} \\
f & \shortstack[l]{Detached lanthanides\\and actinides below} & \shortstack[l]{Lanthanides,\\actinides} & \shortstack[l]{Very similar within rows;\\many radioactive\\(actinides)} \\
\bottomrule
\end{tabular}
\end{table}

The s and p blocks together form the \textbf{main-group elements}. Their behavior depends mainly on outer-electron count, producing the neatest patterns and most of the chemistry encountered in secondary school. The d block contains \textbf{transition elements}, or transition metals. Their differences lie in the next-inner shell while outer shells are nearly alike, producing gradual changes across a row and often several valences (Chapter 15). The f-block lanthanides and actinides are ``exiled'' below solely to keep the table narrow enough to print (Chapter 10).

A practical mnemonic: for a main-group element, the units digit of its group number gives its outer-electron count. Carbon and silicon in group 14 have 4; oxygen and sulfur in group 16 have 6; halogens in group 17 have 7. Outer-electron count gives a rough prediction of valence and behavior.

\section{Common Valences at a Glance}

\textbf{Valence}, also interpretable as oxidation number for ionic compounds, keeps the bonding-electron accounts: surrendered electrons count positive, acquired electrons negative. It is useful bookkeeping, not a literal +2 label pasted onto an atom. Main-group common valences can be read directly from outer-electron counts:

\begin{table}[htbp]
\centering
\begin{tabular}{llll}
\toprule
Group & \shortstack[l]{Outer\\electrons} & \shortstack[l]{Common valences\\(ions)} & Examples \\
\midrule
1 (except H) & 1 & $+1$ & $\chem{Na^+}$, $\chem{K^+}$ \\
2 & 2 & $+2$ & $\chem{Mg^{2+}}$, $\chem{Ca^{2+}}$ \\
13 & 3 & $+3$ & $\chem{Al^{3+}}$ \\
14 & 4 & \shortstack[l]{$-4$ to $+4$; often\\shares electrons} & $\chem{CH_4}$, $\chem{CO_2}$, $\chem{SiO_2}$ \\
15 & 5 & $-3$, $+3$, $+5$ & $\chem{NH_3}$, $\chem{NO_3^-}$, $\chem{PO_4^{3-}}$ \\
16 & 6 & $-2$, $+4$, $+6$ & $\chem{H_2O}$, $\chem{SO_4^{2-}}$ \\
17 & 7 & Mainly $-1$ & $\chem{Cl^-}$, $\chem{NaCl}$ \\
18 & 8 (He: 2) & \shortstack[l]{Usually 0;\\rarely bonds} & $\chem{He}$, $\chem{Ne}$ \\
\bottomrule
\end{tabular}
\end{table}

The rule is: \textbf{a main-group element's highest positive valence usually equals its group's units digit, and its lowest negative valence equals that digit minus 8}. Group 16 oxygen and sulfur, for example, have maximum $+6$, as sulfur in sulfate, and minimum $-2$, as $\chem{O^{2-}}$ and $\chem{S^{2-}}$. The number $8$ is a filled-floor capacity: losing all outer electrons exposes a filled inner floor, while acquiring enough fills the outer floor. These routes correspond to two valences. Atoms do not react with the goal of reaching eight; these arrangements simply happen to have lower energy, as Chapter 10's misconception box explains.

Transition metals resist this mnemonic: iron has $+2$ and $+3$, copper $+1$ and $+2$, and manganese ranges from $+2$ to $+7$. Some next-inner d electrons can participate in bonding; see Chapter 15. The f block most commonly has $+3$, with details beyond this book.

Two details can trip you up. Hydrogen occupies group 1 but is usually not a metal. With nonmetals it generally shares electrons and counts as $+1$, as in water and methane; with very reactive metals it can gain an electron and count as $-1$, as in sodium hydride, $\chem{NaH}$. It is the table's least conforming square; force no group's mnemonic onto it. Second, naming: when one element has several positive valences, Chinese uses prefixes distinguishing higher and lower states. $\chem{Fe^{3+}}$ is iron(III), or ferric, and $\chem{Fe^{2+}}$ iron(II), or ferrous. Iron supplements and rust often contain these two different forms.

\section{Atomic Radius: Smaller Right, Larger Down}

Read the table as a trend map, described here in words. Across a period, atomic radius \textbf{generally decreases}; down a group, it \textbf{increases markedly in steps}. Broadly, upper-right nonmetal atoms are smallest, while lower-left cesium and francium are largest.

\textbf{The reason} is Chapter 11's tug-of-war. Across a period, electrons enter the same floor, but nuclear protons increase, strengthening positive charge and pulling that shell inward. Radius shrinks. Down a group, each row adds another floor: outer electrons are farther away, with full inner shells shielding much of the nuclear attraction like thick blankets. Radius grows. In short, \textbf{right means the same floor with a stronger nuclear pull; down means more floors}.

Real magnitudes make this concrete. Across period 2 from lithium to fluorine, covalent radius falls from approximately \ud{130}{pm} to \ud{60}{pm}, more than halving across seven elements. Down group 1 from lithium to cesium, metallic radii rise from approximately \ud{150}{pm} to \ud{270}{pm}. A picometer, pm, is $10^{-12}$ meters; all atomic radii are on the order of $10^{-10}$ meters. A hundred million atoms in a row are approximately as long as a grain of rice. Trends describe differences of tens of percent on this scale---seemingly small, yet enough to distinguish sodium exploding in water from neon doing nothing.

\textbf{Exceptions and details}: first, contraction across d and f blocks is particularly gradual. Ten transition metals in one row have nearly similar radii; the f block exhibits lanthanide contraction, discussed in the challenge box. Second, downward increments decrease: lithium to sodium differs greatly, rubidium to cesium less. Third, noble-gas radii cannot be compared directly with their neighbors; see the scope below.

\textbf{Scope}: an atomic radius is not a fixed marble diameter. Atoms have probabilistic cloud boundaries, so different measurements yield different values: metallic radii in crystals, covalent radii in molecules, and van der Waals radii in dilute gases, usually the largest. Many tables can give only van der Waals radii for noble gases, making them appear larger than left-hand neighbors. That is a difference of definitions, not an exception. Compare only the same kind of radius, focusing on trends rather than decimal precision.

\section{Ionization Energy: Harder Right, Easier Down}

\textbf{Ionization energy} is the minimum energy required to remove an electron from a gaseous atom. The commonly used first ionization energy is the cost of removing its outermost electron. Across a period, it \textbf{rises unevenly}---notice unevenly---while down a group it \textbf{falls step by step}. Upper-right noble gases are hardest to ionize; lower-left cesium is easiest.

\textbf{The reason} is the other side of the radius coin. Small radius places outer electrons closer to the nucleus under stronger effective attraction, making removal harder. Thus ionization energy rises where radius falls. Down a group, more floors and shielding weaken the hold on outer electrons, lowering the cost. This also explains increasingly vigorous group 1 reactions with water: potassium's outer electron is easier to relinquish than sodium's (Chapter 12).

\textbf{Exceptions}: short periods have two famous small steps. Group 2 is higher than group 13: beryllium > boron and magnesium > aluminum. Group 15 is higher than group 16: nitrogen > oxygen and phosphorus > sulfur. Intuitively, boron loses the lone electron newly installed in a p room, while beryllium must break up a paired set of s electrons, a more stable arrangement. Oxygen has the opposite problem: a pair crowds one p room and repels internally, so removing one relieves crowding. The trend is a sawtooth rise, not a smooth slope.

\textbf{Scope}: ionization energy is defined for a \textbf{single gaseous atom}. It helps explain broad tendencies such as sodium losing and chlorine gaining electrons, but does not directly predict a reaction's speed or heat release. Those also depend on bond-formation energy, solvent, concentration, and many other factors (Chapters 27 and 28). Low ionization energy does not mean reacting with absolutely everything.

The trend has striking applications. Cesium, with the lowest ionization energy, can surrender electrons when illuminated; older photocells and light sensors exploit this light-induced discharge. Conversely, helium and neon have such high ionization energies that bombardment struggles to remove electrons. In neon lamps, electron impacts merely excite them into glowing, leaving them intact afterward, allowing lamps to last more than a decade. The same trend gives opposite corners of the table different livelihoods.

\begin{zhengju}
Ionization energies contain some of the most direct evidence of electron shells. Remove electrons successively from the same atom: the first costs energy, the second more because the remaining positive ion holds tighter, and so on. Sodium, with 11 electrons, has a first ionization energy of approximately \ud{496}{kJ/mol}, then a second of approximately \ud{4560}{kJ/mol}, nearly nine times greater. The third through ninth rise more gradually, before another steep jump at the tenth.

These numbers are a cross-section of the floors. Sodium's one outer electron is easily removed. The next occupies a filled lower shell, immediately raising cost by an order of magnitude. The next eight electrons belong to one floor and show gradual increases; the final electron occupies the innermost first floor and costs most. Other elements behave similarly: magnesium jumps after the second removal and aluminum after the third. \textbf{The position of the successive-ionization-energy step equals the outer-electron count}. Early twentieth-century physicists used such spectra and ionization data to establish that electrons occupy discrete, fixed-capacity floors rather than a shapeless soup.
\end{zhengju}

\section{Electronegativity: Greediest at Upper Right}

\textbf{Electronegativity} measures how strongly an atom pulls shared electrons toward itself \textbf{when bonded to another atom}. Of the three trends, it alone makes sense only in a relationship. It increases left to right and decreases downward, making upper-right fluorine the champion, assigned 4.0 on the Pauling scale, and lower-left elements more generous, with cesium approximately 0.8.

\textbf{The reason} is again the tug-of-war, now in shared accommodation. Strong nuclear attraction and small radius pull the shared pair closer. Fluorine excels at both; moving left or down weakens attraction or increases distance, reducing the pull.

\textbf{Exceptions}: noble gases are generally assigned no electronegativity. This does not mean zero. The property asks who pulls harder while bonding, and noble gases do not bond, making the question inapplicable. The few that can be forced into bonds, such as xenon, have special scales, a later topic.

Consider three everyday judgments. Hydrogen and oxygen have an intermediate electronegativity difference, pulling water's electrons toward oxygen and making it polar. Water can dissolve salt but not oil, whose carbon--hydrogen bonds are almost evenly shared and nonpolar; Chapter 21 develops like-dissolves-like. Fluoride toothpaste exploits fluorine's extreme electronegativity: it grabs electrons more strongly than chlorine and holds tighter after entering tooth hydroxyapatite, making enamel more acid-resistant. Chlorine disinfection and iodine tincture likewise exploit group 17's electron-grabbing oxidizing power. The same family greed can help or harm depending on dose (Chapter 13).

\textbf{Scope}: remember three points. First, electronegativity is a \textbf{scale value}, not a directly weighed quantity. Pauling's scale is most common; others differ, so compare within one scale. Second, it predicts \textbf{which way a covalent bond's electrons lean}. A large difference, as for sodium and chlorine, nearly transfers the electron outright, giving an ionic bond. A small difference, as for carbon and hydrogen, shares nearly evenly. An intermediate one, as for hydrogen and oxygen, gives a polar covalent bond leaning toward the more electronegative atom. Water's partly positive hydrogen ends and partly negative oxygen end follow from this, along with many unusual properties (Chapter 21). Third, the difference is an empirical guide: bonding is a continuum without a sharp dividing line, and electronegativity does not predict reaction rates.

\begin{qiaoliang}
The three trends are three sides of one principle, condensed into two statements:

\textbf{Across}: the same floor, more protons, unchanged shielding, stronger effective attraction---smaller radius, higher ionization energy, higher electronegativity.

\textbf{Down}: more floors and shielding, outer electrons farther away and screened---larger radius, lower ionization energy, lower electronegativity.

The table therefore has an overall diagonal trend: toward the upper right, excluding noble gases, atoms are smaller, hold electrons tighter, and are greedier; toward the lower left, atoms are larger, electrons looser, and generosity greater. Metal--nonmetal boundaries and family reactivity rankings follow from this. Instead of memorizing six arrows, remember the tug-of-war and changing floor count and derive them on the spot.
\end{qiaoliang}

\begin{wuqu}
``The more electronegative atom carries a full negative charge in a molecule.'' No. It merely \textbf{pulls the shared pair closer}, without complete transfer. Water's oxygen is partially negative, often $\delta^-$, corresponding to a fraction of an electron's charge, far from a full $\chem{O^{2-}}$. Only combinations with extremely large differences and genuinely transferred electrons, such as $\chem{NaCl}$, have full ionic charges. Confusing partial negative charge with an anion mispredicts many properties, including why water does not readily ionize completely (Chapters 21 and 25).
\end{wuqu}

\section{Meet an Element in Thirty Seconds}

The manual's real value is letting you sketch an unfamiliar element immediately. Try strontium, \chem{Sr}, number 38:

\begin{itemize}[itemsep=2pt]
\item \textbf{Locate}: period 5, group 2, s block, between calcium above and barium below.
\item \textbf{Guess valence}: group 2 means two outer electrons and almost always $+2$, giving $\chem{Sr^{2+}}$. Copy calcium's compound patterns: strontium chloride $\chem{SrCl_2}$ and carbonate $\chem{SrCO_3}$.
\item \textbf{Guess behavior}: ionization energy falls and metallic reactivity rises down the group. Strontium should lose electrons more eagerly and react with water more vigorously than calcium, requiring isolation from air like sodium. Its flame color supplies red fireworks, its main everyday appearance.
\end{itemize}

Now try the right side: selenium, \chem{Se}, number 34, in period 4, group 16, below oxygen and sulfur. Guess valences from $-2$ to $+6$ and copy sulfur's compound patterns. Its electronegativity is slightly below sulfur's, so it grabs electrons somewhat less strongly. It is an essential trace nutrient with a narrow safe range: excess poisons, deficiency causes illness, explaining Chapter 14's repeated emphasis on dose. In thirty seconds, an unfamiliar name becomes an acquaintance with coordinates, valences, and temperament.

This is both the table's limit and its dignity as a tool: it always provides \textbf{an informed first guess}, never the final answer. Experiments decide whether the guess is right. Before germanium was discovered, Mendeleev's eka-silicon was also a guess (Chapter 10).

\begin{shiyan}
\textbf{Periodic-table treasure hunt}: just a printed periodic table, paper, and pen; no chemicals, so safe.

Materials: a periodic table with atomic numbers, from the back of the book or printed online, plus paper and pen.

Procedure: without reading explanatory text, answer from coordinates alone, then check against an annotated table:
\begin{enumerate}[itemsep=1pt]
\item Find the element in period 4, group 16, write its symbol, and predict its lowest negative valence.
\item Between potassium, \chem{K}, and rubidium, \chem{Rb}, guess which reacts more vigorously with water and explain in one sentence.
\item Locate the most electronegative element and the stable element with strongest metallic character. Are they at the extreme upper right and lower left?
\item Choose an unfamiliar transition metal, find the element below it, and research their uses. Does family resemblance hold in the d block?
\end{enumerate}

What to observe: how many predictions were correct? For incorrect ones, was the trend inapplicable, as for noble gases, or did you skip a reasoning step? Write the cause beside each mistake; this teaches more than getting everything right.
\end{shiyan}

\begin{tiaozhan}
\textbf{Lanthanide contraction and diagonal similarities}: optional; you can use the tool without this.

Period 6 transition metals such as hafnium, tantalum, and tungsten have nearly the same radii as their period 5 counterparts, zirconium, niobium, and molybdenum, and can be so similar that separation is difficult. The reason hides in the f block: fourteen lanthanide elements add electrons to deep 4f rooms, whose shielding is weak. Effective nuclear attraction quietly increases across the series, shrinking the subsequent period 6 atoms. This lanthanide contraction makes the usual downward-radius increase nearly disappear in the d block.

Another small pattern is the diagonal relationship: short-period diagonal neighbors lithium--magnesium, beryllium--aluminum, and boron--silicon are unexpectedly similar because rightward shrinkage and downward growth partly cancel. These are not facts to memorize, but reminders that periodic behavior combines competing factors. At the table's edges, their combination produces variations: the rules still apply, but require finer accounting.
\end{tiaozhan}

\section{What This Manual Cannot Do}

Finally, the boundaries. Trends and valence tables are \textbf{a first-guess production line}, not a destination. They cannot predict reaction rates, the kinetics of Chapter 27, or final reaction extent, the equilibria and free energy of Chapter 28, much less molecular three-dimensional shapes, covered by Appendix C and Chapters 19 and 20. For a particular substance, use these guesses as starting questions and keep digging with the five recurring questions.

When you want to know how elements assemble into molecules and molecules into life, turn to Appendix C, Reading Molecular Models, the next tool. If a symbol, formula, or arrow itself blocks you, Appendix A, A First-Aid Kit for Chemical Symbols, handles that level. Their roles are distinct: this appendix predicts an element's broad behavior and valences; Appendix C explains the shapes of bonded molecules and what different drawings preserve.

\section*{Exercises}

\begin{sikaoti}
\item Without consulting a table, compare chlorine, period 3 group 17, and bromine, period 4 group 17. Which has larger radius, higher first ionization energy, and higher electronegativity? Give one sentence of reasoning for each.
\item Aluminum's successive ionization energies jump sharply between the third and fourth removals. Without looking up data, sketch the expected staircase with removal number horizontally and energy vertically. What does the jump's position show?
\item A student writes, ``Oxygen is more electronegative than hydrogen, so oxygen in water has two negative charges and each hydrogen one positive charge.'' Identify the mistaken level of reasoning and rewrite correctly.
\item Apply the thirty-second procedure to rubidium, \chem{Rb}, number 37, and iodine, \chem{I}, number 53. Locate each, predict valences and one physical or chemical behavior, then research and record errors.
\item Why must the radius trend be qualified for noble gases? How would you label their radii on a more honest wall chart to prevent misunderstanding?
\item Gallium, \chem{Ga}, lies below aluminum and above and slightly right of zinc. Combining block, radius, and ionization trends, predict its highest positive valence and whether it is a metal. Then look up its melting point: it can melt in your palm. Could that property be inferred directly from periodic trends, and why?
\end{sikaoti}

\textbf{Translator safety note:} The water-reaction, flame-color, and palm-melting examples are descriptions or research questions, not handling instructions. Do not test reactive metals, make fireworks, or put chemical samples on your skin. Use published observations and follow \href{https://institute.acs.org/acs-center/lab-safety/education-training/safer-experiments.html}{ACS experiment risk-assessment guidance} for supervised activities.
